\section{Quantizing the Dirac Field}

\subsection{A Glimpse at the Spin--Statistics Theorem (What You Need)}
Empirical + theoretical result:
\begin{equation}
\text{integer spin} \Rightarrow \text{bosons (commutators)}, 
\qquad
\text{half-integer spin} \Rightarrow \text{fermions (anticommutators)}.
\end{equation}

Why it matters (minimal statement):
\begin{itemize}
\item Using commutators for spin-$\frac12$ fields leads to negative-norm states or violation of causality.
\item Using anticommutators ensures positivity of energy and consistent relativistic causality.
\end{itemize}

\subsection{Canonical Fermionic Quantization}
Start from the Dirac Lagrangian:
\begin{equation}
\mathcal{L}_D=\bar{\psi}(i\gamma^\mu\partial_\mu-m)\psi.
\end{equation}
Canonical momentum conjugate to $\psi$:
\begin{equation}
\pi_\psi := \frac{\partial\mathcal{L}}{\partial(\partial_0\psi)} = i\psi^\dagger.
\end{equation}

Impose equal-time \textbf{anticommutation relations}:
\begin{equation}
\{\psi_\alpha(\mathbf{x},t),\psi_\beta^\dagger(\mathbf{y},t)\}
=
\delta_{\alpha\beta}\,\delta^{(3)}(\mathbf{x}-\mathbf{y}),
\end{equation}
\begin{equation}
\{\psi_\alpha(\mathbf{x},t),\psi_\beta(\mathbf{y},t)\}=0,
\qquad
\{\psi_\alpha^\dagger(\mathbf{x},t),\psi_\beta^\dagger(\mathbf{y},t)\}=0.
\end{equation}
Here $\alpha,\beta$ are spinor indices ($1,\dots,4$).

\subsection{Mode Expansion of the Dirac Field}
The free Dirac field operator is expanded as
\begin{equation}
\psi(x)
=
\int \frac{d^3p}{(2\pi)^3}\frac{1}{\sqrt{2E_{\mathbf{p}}}}
\sum_{s=1,2}
\left(
b_{\mathbf{p}}^{(s)}u^{(s)}(p)e^{-ip\cdot x}
+
d_{\mathbf{p}}^{(s)\dagger}v^{(s)}(p)e^{ip\cdot x}
\right),
\qquad p^0=E_{\mathbf{p}}.
\end{equation}
Interpretation:
\begin{itemize}
\item $b_{\mathbf{p}}^{(s)}$ annihilates a particle (fermion).
\item $d_{\mathbf{p}}^{(s)}$ annihilates an antiparticle (antifermion).
\end{itemize}

Anticommutators:
\begin{equation}
\{b_{\mathbf{p}}^{(r)},b_{\mathbf{q}}^{(s)\dagger}\}
=
(2\pi)^3\delta^{(3)}(\mathbf{p}-\mathbf{q})\delta^{rs},
\qquad
\{d_{\mathbf{p}}^{(r)},d_{\mathbf{q}}^{(s)\dagger}\}
=
(2\pi)^3\delta^{(3)}(\mathbf{p}-\mathbf{q})\delta^{rs},
\end{equation}
all others vanish.

\subsection{Fermi--Dirac Statistics and the Pauli Principle}
Because operators anticommute,
\begin{equation}
\left(b_{\mathbf{p}}^{(s)\dagger}\right)^2=0,
\qquad
\left(d_{\mathbf{p}}^{(s)\dagger}\right)^2=0,
\end{equation}
so no two identical fermions can occupy the same one-particle state:
\begin{equation}
n_{\mathbf{p}}^{(s)}\in\{0,1\}.
\end{equation}
This is the Pauli exclusion principle.

Number operators:
\begin{equation}
N_b=\int\frac{d^3p}{(2\pi)^3}\sum_s b_{\mathbf{p}}^{(s)\dagger}b_{\mathbf{p}}^{(s)},
\qquad
N_d=\int\frac{d^3p}{(2\pi)^3}\sum_s d_{\mathbf{p}}^{(s)\dagger}d_{\mathbf{p}}^{(s)}.
\end{equation}

\subsection{Hamiltonian and Energy Positivity}
The free Dirac Hamiltonian can be written (after normal ordering) as
\begin{equation}
:H:
=
\int\frac{d^3p}{(2\pi)^3}\sum_s E_{\mathbf{p}}
\left(
b_{\mathbf{p}}^{(s)\dagger}b_{\mathbf{p}}^{(s)}
+
d_{\mathbf{p}}^{(s)\dagger}d_{\mathbf{p}}^{(s)}
\right).
\end{equation}
Hence both particles and antiparticles contribute positive energy.

\subsection{Particles and Antiparticles}
Define vacuum:
\begin{equation}
b_{\mathbf{p}}^{(s)}|0\rangle=0,
\qquad
d_{\mathbf{p}}^{(s)}|0\rangle=0
\quad \forall\,\mathbf{p},s.
\end{equation}
One-particle states:
\begin{equation}
|\mathbf{p},s\rangle = b_{\mathbf{p}}^{(s)\dagger}|0\rangle,
\qquad
|\overline{\mathbf{p},s}\rangle = d_{\mathbf{p}}^{(s)\dagger}|0\rangle.
\end{equation}

Conserved $U(1)$ current:
\begin{equation}
j^\mu=\bar{\psi}\gamma^\mu\psi,
\qquad
Q=\int d^3x\,j^0.
\end{equation}
In terms of operators (schematically):
\begin{equation}
Q=\int\frac{d^3p}{(2\pi)^3}\sum_s
\left(
b_{\mathbf{p}}^{(s)\dagger}b_{\mathbf{p}}^{(s)}
-
d_{\mathbf{p}}^{(s)\dagger}d_{\mathbf{p}}^{(s)}
\right).
\end{equation}
Thus antiparticles carry opposite charge.

\subsection{Dirac's Hole Interpretation (Historical Motivation)}
The Dirac equation has negative-energy solutions. Dirac proposed:
\begin{itemize}
\item All negative-energy states are filled in the vacuum (Dirac sea).
\item A missing negative-energy electron appears as a positively charged particle (the positron).
\end{itemize}
Modern QFT interpretation: negative-energy modes are reinterpreted as antiparticles with positive energy.

\subsection{Fermion Propagator}
The \textbf{Feynman propagator} for the Dirac field:
\begin{equation}
S_F(x-y)
:=
\langle 0|T\{\psi(x)\bar{\psi}(y)\}|0\rangle.
\end{equation}
Momentum-space representation:
\begin{equation}
S_F(x-y)
=
\int\frac{d^4p}{(2\pi)^4}\,
\frac{i(\slashed{p}+m)}{p^2-m^2+i\epsilon}\,e^{-ip\cdot(x-y)}.
\end{equation}
It satisfies the Green function equation:
\begin{equation}
(i\slashed{\partial}_x-m)\,S_F(x-y)=i\delta^{(4)}(x-y).
\end{equation}

\subsection{Spin Sums (Used Constantly)}
\begin{equation}
\sum_s u^{(s)}(p)\bar{u}^{(s)}(p)=\slashed{p}+m,
\qquad
\sum_s v^{(s)}(p)\bar{v}^{(s)}(p)=\slashed{p}-m.
\end{equation}

\subsection{Wick's Theorem for Fermions (Sign Rules)}
Wick's theorem also holds for fermions, but each exchange of fermionic operators introduces a minus sign.

Fermion contraction:
\begin{equation}
\contraction{}{\psi(x)}{}{\bar{\psi}(y)}\psi(x)\bar{\psi}(y)
=
\langle 0|T\{\psi(x)\bar{\psi}(y)\}|0\rangle
=
S_F(x-y).
\end{equation}
Time ordering for fermions includes signs:
\begin{equation}
T\{\psi(x)\bar{\psi}(y)\}
=
\theta(x^0-y^0)\psi(x)\bar{\psi}(y)
-\theta(y^0-x^0)\bar{\psi}(y)\psi(x).
\end{equation}

\subsection{Feynman Rules (Dirac Fermions, Minimal Set)}
For a theory with interaction $\mathcal{L}_{\mathrm{int}}$:
\begin{equation}
S=T\exp\left(i\int d^4x\,\mathcal{L}_{\mathrm{int}}(x)\right).
\end{equation}

\paragraph{Free fermion line (propagator).}
Internal fermion line with momentum $p$:
\begin{equation}
\frac{i(\slashed{p}+m)}{p^2-m^2+i\epsilon}.
\end{equation}

\paragraph{External spinors.}
For external fermion legs:
\begin{equation}
\text{incoming fermion: } u^{(s)}(p),
\qquad
\text{outgoing fermion: } \bar{u}^{(s)}(p),
\end{equation}
\begin{equation}
\text{incoming antifermion: } \bar{v}^{(s)}(p),
\qquad
\text{outgoing antifermion: } v^{(s)}(p).
\end{equation}

\paragraph{Closed fermion loops.}
Each closed fermion loop contributes an extra factor:
\begin{equation}
(-1).
\end{equation}

\paragraph{Example: Yukawa vertex.}
If
\begin{equation}
\mathcal{L}_{\mathrm{int}}=-g\,\bar{\psi}\psi\phi,
\end{equation}
then the vertex factor is
\begin{equation}
-i g.
\end{equation}

\paragraph{Example: QED vertex.}
If
\begin{equation}
\mathcal{L}_{\mathrm{int}}=-e\,\bar{\psi}\gamma^\mu\psi A_\mu,
\end{equation}
then vertex factor:
\begin{equation}
-ie\gamma^\mu.
\end{equation}

\subsection{Summary (Exam Checklist)}
\begin{itemize}
\item Fermions must be quantized with anticommutators:
\[
\{\psi_\alpha(\mathbf{x}),\psi_\beta^\dagger(\mathbf{y})\}=\delta_{\alpha\beta}\delta^{(3)}(\mathbf{x}-\mathbf{y}).
\]
\item Mode expansion introduces particle operators $b$ and antiparticle operators $d$.
\item Pauli principle: $(b^\dagger)^2=0$ and occupancy $n\in\{0,1\}$.
\item Dirac propagator:
\[
S_F(p)=\frac{i(\slashed{p}+m)}{p^2-m^2+i\epsilon}.
\]
\item Spin sums:
\[
\sum u\bar{u}=\slashed{p}+m,\quad \sum v\bar{v}=\slashed{p}-m.
\]
\item Closed fermion loop $\Rightarrow$ extra factor $(-1)$.
\end{itemize}
